\addcontentsline{toc}{subsubsection}{Solution 48. The One-Point Cubic Lock}

\begin{dlsummary}[title={\bfseries Solution 48. The One-Point Cubic Lock},breakable]

\begin{proof}
The unique solution is
\[
\boxed{f(x)=x^3\qquad(x>0).}
\]

\paragraph{1. Extract a global cubic lower bound.}

Set
\[
g(x)=f(x)-x^3.
\]
Because
\[
(x+y)^3=x^3+y^3+3xy(x+y),
\]
the first inequality becomes
\[
g(x+y)\ge g(x)+g(y). \tag{1}
\]

For any positive integer \(n\), repeated application gives
\[
g(t)\ge n\,g(t/n).
\]
Returning to \(f\), and using its positivity,
\[
f(t)\ge n f(t/n)+t^3\left(1-\frac1{n^2}\right)
>t^3\left(1-\frac1{n^2}\right).
\]
Letting \(n\to\infty\), we obtain
\[
\boxed{f(t)\ge t^3,\quad\text{so }g(t)\ge0.} \tag{2}
\]

\paragraph{2. Lock the value at \(2\).}

Applying (2) twice,
\[
512\ge f(f(2))\ge f(2)^3\ge8^3=512.
\]
Therefore every inequality is an equality, and
\[
\boxed{f(2)=8,\qquad g(2)=0.} \tag{3}
\]

\paragraph{3. Determine the function on an interval.}

For \(0<t<2\), superadditivity gives
\[
0=g(2)\ge g(t)+g(2-t).
\]
Both summands are nonnegative, so \(g(t)=0\). Hence
\[
\boxed{f(t)=t^3\qquad(0<t\le2).} \tag{4}
\]

\paragraph{4. Propagate to every positive input.}

Fix \(x>0\). Choose a positive integer \(k\) such that \(x/2^k<2\). Repeated submultiplicativity, followed by (3) and (4), gives
\[
f(x)
\le f(2)^k f(x/2^k)
=8^k\left(\frac{x}{2^k}\right)^3
=x^3.
\]
Combined with (2), this proves \(f(x)=x^3\).

\textbf{Verification.} The cubic identity makes the first inequality an equality; multiplicativity also holds, and
\[
f(f(2))=f(8)=512.
\]
\end{proof}

\paragraph{Takeaways.}
\begin{itemize}
\item **Subtract the expected model:** Removing \(x^3\) exposes an ordinary superadditive defect.
\item **Subdivide before taking limits:** Positivity of \(f\) controls the potentially negative defect at small inputs.
\item **One exact value can determine an interval:** A nonnegative superadditive function that vanishes at \(2\) must vanish throughout \((0,2)\).
\item **Use different inequalities locally and globally:** Addition determines a bounded interval; multiplication propagates the result everywhere.
\end{itemize}

\end{dlsummary}
